\section{結論}
\label{sec:conclusion}

本研究では，パターンサポート時系列に合成コントロール法を適用し，サポート変動を外部イベントに因果帰属するための初の手法 SC-DenseAttrib を提案した．
本フレームワークは，SUTVA と線形因子モデルを含む因果識別の仮定を明示化し，ドナープール，反事実サポート軌跡，サポートに対する因果効果，およびアイテム非共有パターン制御を形式的に定義した．
線形因子モデルのもとでの$\sqrt{t_0}$-一致性，およびプラセボ検定の正確なp値を理論的に保証した．

7つの実験により，SC-DenseAttrib が以下を達成することが示された：
(1) 平均絶対誤差 0.04 未満で因果効果を正確に回復する；
(2) DiD，Before-After，アイテム共有 SCM の各ベースラインを上回る推定精度を示す；
(3) 合成コントロール構成により季節性の交絡を処理する；
(4) SUTVA 違反が真の効果の5\%以下であれば推定精度の劣化が軽微であることを確認した．

\paragraph{制約}
主な制約は以下の通りである：
(1) 統計的検出力がドナープールサイズに依存し，$p < 0.05$ には $J \geq 19$ が必要である；
(2) SUTVA は完全には保証されず，間接的波及効果が存在する場合にバイアスが発生する；
(3) 本研究の実験は合成データに限定されており，実世界データでの検証が今後必要である；
(4) ドナー数が大きい場合の過適合リスクに対して，正則化（\Cref{rem:regularization}）の導入が望ましい．

\paragraph{今後の課題}
(1) Augmented SCM~\cite{benmichael2021augmented} やリッジ正則化による推定量の改善；
(2) 条件付き独立性検定によるアイテム非共有基準の緩和；
(3) 密集区間検出パイプラインとの統合による完全自動化された因果帰属；
(4) Dunnhumby 等の実世界小売データセットへの適用と外的妥当性の検証．
